\chapter{Series y funciones}
\label{cap:series-funciones}

Las series de funciones combinan las ideas de convergencia de series numéricas con el comportamiento de funciones. Prestaremos atención a la convergencia puntual y uniforme, y a criterios que permiten demostrar convergencia uniforme de manera práctica \cite{abbott2015,rudin1976}.

\section{Series de funciones}
Convergencia puntual y uniforme de series y convergencia normal.

\section{Criterio \texorpdfstring{$M$}{M} de Weierstrass}
Prueba de convergencia uniforme a partir de cotas sumables independientes del punto.

\section{Series de potencias}
Radio e intervalo real de convergencia y operaciones válidas en su interior.

\section{Series de Taylor y polinomios de Taylor}
Series asociadas a las derivadas y condiciones para que la serie de Taylor coincida con la función.

\section{Ejemplos y aplicaciones}
Aplicaciones de las series de funciones a aproximaciones y al estudio de funciones.

\subsection{La irracionalidad de \texorpdfstring{$e$}{e} y de \texorpdfstring{$\pi$}{pi}}
Aplicaciones de desarrollos en series para estudiar la irracionalidad de constantes clásicas.

\subsection{Una función continua que no es derivable en ningún punto}
Construcción de una función de Weierstrass mediante una serie uniformemente convergente.

\subsection{Teorema de aproximación de Weierstrass}
Aproximación uniforme de funciones continuas mediante polinomios.
